\documentclass[10pt,a4paper]{amsart}
\usepackage[margin=19mm]{geometry}
\usepackage{amsmath,amssymb,mathtools}
\usepackage[scr=rsfs,cal=euler]{mathalfa}
\usepackage{xfrac}
\usepackage[unicode,hidelinks,bookmarks=false]{hyperref}
\newcommand{\ie}{i.e.\ }
\newcommand{\cf}{cf.\ }
\newcommand{\resp}{respectively\ }
\newcommand{\Subsection}{\subsection}
\providecommand{\nobreakdash}{\nobreak\hbox{-}}
\setlength{\emergencystretch}{2em}
\multlinegap0pt
\usepackage[shortlabels]{enumitem}
\usepackage{bm}
\usepackage[normalem]{ulem}
\usepackage{amssymb}
\usepackage{mathtools}
\usepackage{tikz}
\newtheorem{lemma}{Lemma}
\newtheorem{corollary}{Corollary}
\newtheorem{theorem}{Theorem}
\usepackage{cleveref}
\crefname{lemma}{Lemma}{Lemmas}
\crefname{corollary}{Corollary}{Corollarys}
\crefname{theorem}{Theorem}{Theorems}
\newcommand{\E}{\textsf{\upshape E}}
\newcommand{\Gnp}{\mathcal{G}(n,p)}
\newcommand{\prob}{\textsf{\upshape Pr}}
\title{Multiset Metric Dimension of Binomial Random Graphs}
\author{Austin Eide, Pawel Pralat}
\date{Source: arXiv:2507.11686v2 (CC BY 4.0)}
\begin{document}
\maketitle

\noindent\emph{Source and licence.} Multiset Metric Dimension of Binomial Random Graphs. Version arXiv:2507.11686v2 (CC BY 4.0), distributed by arXiv under the Creative Commons Attribution 4.0 International licence, http://creativecommons.org/licenses/by/4.0/, as recorded in the arXiv licence metadata for this exact version and shown on the abstract page https://arxiv.org/abs/2507.11686v2. Full licence notice: \texttt{licenses/arxiv-2507.11686-CC-BY-4.0.txt}.\par\smallskip

\noindent\textit{Editorial scope.} Counted unit: the article's theorem thm:upper\_bound (source0.tex lines 388--389). The article's complete original proof of it is delivered with it (source0.tex lines 402--457, one \texttt{proof} environment of 56 lines, which the article places away from the statement). No mathematics is written, repaired or re-derived: the statement and the proof are the article's own lines, in the article's own notation, and no retained line is substituted or re-flowed.  Retained uncounted as the source-local context the proof needs: lines 255--260; lines 347--355; lines 393--398.  Nothing was omitted from any retained span, so the package carries no omission record.  No retained line contains a citation, so the wrapper carries no bibliography and no bibliography entry.  No retained line contains a figure environment, an image-inclusion command or drawing code, so no image is delivered and the package declares no asset dependency.  Editorial formatting only, in addition: an AMS \texttt{amsart} preamble that restates the article's own declarations and macros used by the retained lines, namely the environments \texttt{lemma}, \texttt{corollary}, \texttt{theorem} on separate counters, together with the standard packages those lines need.  The retained environments print in this document as Lemma 1, Corollary 1, Theorem 1 in order of appearance, whereas the article prints the same items with the numbering of its own full document.  The complete native source of the article as distributed in the arXiv e-print for this exact version is bundled unchanged as \texttt{sources/181645/source0.tex}.
\begin{lemma}\label{lem:exponent_func2}
	Let $x \in (0,1]$ and $y \in [0,1]$ be such that $f_{x}(y) > 0$ (or, equivalently, that $1-\left\lfloor\frac{1}{x}\right\rfloor x < y \leq 1$). Then for any $\epsilon > 0$, we have
	$$
		f_{x}(y+\epsilon) - f_{x}(y) \geq \epsilon.
	$$
\end{lemma}

\begin{corollary}\label{cor:expansion_cor}
Let $d = (n-1)p$ satisfy $d = \omega(\log n)$ and $d = o(n)$. Then the following hold w.h.p.\ for $\Gnp$: 
	\begin{enumerate}[label = (\roman*)]
		\item For any vertex $v$ and any $i \in \{ 0,1,\dots,i^{*} \} $, $|N_{i}(v)| = (1+o(1))d^{i}$.
		\item For any pair of vertices $v$ and $w$ and any $i  \in \{ 0,1,\dots,i^{*} \} $, $|S_{i}(w) \setminus N_{i}(v)| = (1+o(1))d^{i}$.
		\item If $c = \Theta(1)$, for any vertex $v$, $|N_{i^{*}+1}(v)| = (1+o(1))(1-e^{-c})n$.
		\item If $c = \Theta(1)$, for any pair of vertices $v$ and $w$, $|S_{i^{*}+1}(v) \setminus N_{i^{*}}(w)| = (1+o(1))(1-e^{-c})n$ and $|S_{i^{*}+1}(w) \setminus N_{i^{*}+1}(v)| = (1+o(1))e^{-c}(1-e^{-c})n$.
	\end{enumerate}
\end{corollary}

\begin{lemma}\label{lem:binomial}
	Let $Z$ be a binomial random variable with mean $\mu$. Then 
	$$
		\max_{z}\prob(Z = z) \leq \min\left\{1, O\left(\frac{1}{\sqrt{\mu}} \right)\right\}.
	$$
\end{lemma}

\begin{theorem}\label{thm:upper_bound} Let $x \in \left(0, \frac{1}{8} \right]$ be fixed and let $d = (n-1)p = n^{x + O(\log^{-1}n)}$. W.h.p., $\Gnp$ has a multiset resolving set of size at most $n^{y_{4} +O(\log^{-1}n)}$, where $y_{4} = y_{4}(x) \in (0,1)$ is the unique solution to $f_{x}(y) = 4$.
\end{theorem}

\begin{proof}[Proof of \cref{thm:upper_bound}]
We condition on $\Gnp$ satisfying $\mathcal{V}$. Let $\epsilon = \frac{C}{\log n}$ with $C = C(x) >0$ a large constant to be determined later, let $y = y(n) = y_{4} + \epsilon$, and let $r = r(n) = n^{y} $. Let $R$ be a random subset of $V$ obtained by including each vertex $v$ in $R$ independently with probability $r/n$. Our goal is to show that for each pair $v,w \in V$, 
\begin{equation}\label{eq:to_show}
	\prob(R \text{ fails to distinguish }v\text{ and }w\,|\,\mathcal{V}) = n^{-\frac{1}{2}f_{x}(y) + O(\log^{-1}n)}.
\end{equation}
By Lemma~\ref{lem:exponent_func2} we have $f_{x}(y) \geq f_{x}(y_{4}) + \epsilon = 4 + \frac{C}{\log n}$. Thus we can choose $C$ sufficiently large so that the right-hand side of \eqref{eq:to_show} is at most $\frac{1}{2n^{2}}$, and hence by a union bound we have
$$
	\prob(R\text{ fails to distinguish some pair of vertices}\,|\,\mathcal{V}) \leq \binom{n}{2} \cdot \frac{1}{2n^{2}} \le \frac{1}{4}.
$$
By Markov's inequality the probability that $|R|$ exceeds $2\E[R] = 2r$ is at most $\frac{1}{2}$. Using this and the above, the probability that $R$ distinguishes all vertices and has size at most $2r = n^{y_{4} + O(\log^{-1}n)}$ is at least $\frac{1}{4}$. We may then conclude that any graph satisfying $\mathcal{V}$ has a multiset resolving set of size at most $n^{y_{4} +O(\log^{-1}n)}$, and thus $\Gnp$ has a multiset resolving set of size at most $n^{y_{4} + O(\log^{-1}n)}$ w.h.p.

Now we focus on showing \eqref{eq:to_show}. Fix $v$ and $w$. We will reveal the set $R$ in stages with respect to these two vertices. Initially none of $R$ is revealed. In stage $0$ we reveal $R \cap N_{0}(\{v,w\}) = \{v,w\}$. In stage $1$ we reveal $R \cap (N_{1}(\{v,w\}) \setminus N_{0}(\{v,w\}))$, and in general in stage $i$ we reveal $R \cap (N_{i}(\{v,w\}) \setminus N_{i-1}(\{v,w\}))$. We reveal $R$ in this way through stage $\lfloor 1/x \rfloor$, though we need to take some care to handle the final stage when $x = \frac{1}{k}$ for some integer $k \ge 8$, as in this case we have $i^{*} = \lfloor1/x\rfloor -1$. We start with the analysis of stages $0$ through $i^{*}$, which applies to all $x \in (0,1/8]$. 

In stage $0$, we simply reveal whether $v$ and $w$ are in $R$. We condition on neither vertex being in $r$, as this occurs with probability $1 - O(r/n) = 1-o(1)$. For $i \geq 1$, condition on the outcomes of stages $0,1,\dots,i-1$. To process stage $i$, we first reveal which vertices in $S_{i}(v) \setminus N_{i-1}(w)$ are in $R$ and also condition on these. At this point we have the following information:
	\begin{enumerate}
		\item $|S_{i}^{R}(v)|$ is deterministic (i.e., it is a function of the information revealed so far);
		\item $|S_{i}^{R}(w)|$ can be expressed as $z_{i} + Z_{i}$, where $z_{i} = |S_{i}^{R}(w) \cap N_{i}(v)|$ is deterministic and $Z_{i} = |S_{i}^{R}(w) \setminus N_{i}(v)|$ is a random variable; 
		\item $Z_{i}$ is a $\text{Bin}(s_{i}, r/n)$ random variable with $s_{i} = |S_{i}(w) \setminus N_{i}(v)|$.
	\end{enumerate}
By the event $\mathcal{V}$, we have $s_{i} = (1+o(1))d^{i} = (1+o(1))n^{ix + O(\log^{-1}n)} = \Theta(n^{ix}).$ Observe that $|S_{i}^{R}(v)|=|S_{i}^{R}(w)|$ if and only if $Z_{i} = |S_{i}^{R}(v)| - z_{i}$. Conditioned on the information revealed so far, the probability that $Z_{i} = |S_{i}^{R}(v)| - z_{i}$ is at most
\begin{eqnarray*}
	\max_{z}\prob(\text{Bin}(s_{i}, r/n) = z) &\leq& \min\left\{1, O\left(\sqrt{\frac{n}{rs_{i}}}\right) \right\}\\
	&=&O\left(\min\left\{1, \sqrt{\frac{n}{rn^{ix}}}\right\}\right)\\
	&=&O\left(\min\left\{1, n^{\frac{1-ix-y}{2}} \right\} \right)\\
	&=&O\left(n^{-\frac{1}{2}\max\{ ix+y-1, 0 \} }\right),
\end{eqnarray*}
where we use \cref{lem:binomial} for the inequality in the first line. Suppose $x \in \left(\frac{1}{k+1}, \frac{1}{k} \right)$ for some integer $k \ge 8$. In this case, we have $i^{*} = \lfloor 1/x \rfloor = k$, and from the above we inductively have the bound
\begin{eqnarray*}
	\prob\left(\bm{m}_{R}(v) = \bm{m}_{R}(w)\,|\,\mathcal{V}\right) &\leq& \prob\left(\bigcap_{i=0}^{i^{*}}\left\{|S_{i}^{R}(v)| = |S_{i}^{R}(w)|\right\}\,|\,\mathcal{V}\right)\\
	&\leq& \prod_{i=0}^{i^{*}}O\left(n^{-\frac{1}{2}\max\{ ix+y-1, 0 \} }\right)\\
	&=&O\left(n^{-\frac{1}{2}\sum_{i=0}^{\lfloor 1/x \rfloor}\max\{ix+y-1,0\} }\right)\\
	&=&n^{-\frac{1}{2}f_{x}(y)+ O(\log^{-1}n)},
\end{eqnarray*}
which establishes \eqref{eq:to_show} for this case. 

Suppose now that $x = \frac{1}{k}$ for some $k \ge 8$. For this case $i^{*} = \lfloor 1/x \rfloor -1 = k-1$ and $c = \frac{d^{i^{*}+1}}{n} = n^{O(\log^{-1}n)}= \Theta(1)$. We run stage $i^{*}+1$ in the same way as stages $0$ through $i^{*}$. Let $s_{i^{*}+1} = |S_{i^{*}+1}(w) \setminus N_{i^{*}+1}(v)|$, and observe by \cref{cor:expansion_cor} that $s_{i^{*}+1} = (1+o(1))e^{-c}(1-e^{-c})n = \Theta(n)$. Using \cref{lem:binomial}, we can then upper bound probability that $|S_{i^{*}+1}^{R}(v)| = |S_{i^{*}+1}^{R}(w)|$ (conditioned on the history) by
$$
	\max_{z}\prob(\text{Bin}(s_{i^{*}+1}, r/n) = z) \leq \min\left\{1, O\left(\sqrt{\frac{n}{rs_{i^{*}+1}}}\right) \right\} = O(r^{-1/2}) = n^{-y/2 + O(\log^{-1}n)}.
$$
We then have the bound
\begin{eqnarray*}
	\prob\left(\bm{m}_{R}(v) = \bm{m}_{R}(w)\,|\,\mathcal{V}' \right) &\leq& n^{-\frac{1}{2}\sum_{i=0}^{i^{*}}\max\{ix + y - 1,0\} + O(\log^{-1}n)} \cdot n^{-\frac{1}{2}y + O(\log^{-1}n)}\\
	&=&n^{-\frac{1}{2}\sum_{i=0}^{\lfloor 1/x \rfloor -1}\max\{ix + y - 1,0\} - \frac{1}{2}y + O(\log^{-1}n)}.
\end{eqnarray*}
Note that
\begin{eqnarray*}
	\sum_{i=0}^{\lfloor1/x\rfloor - 1}\max\{ix+y-1,0\} &=& f_{x}(y) - \max\{\lfloor1/x\rfloor x + y-1,0\}\\
	&=& f_{x}(y) - \max\{k \cdot \frac{1}{k} + y - 1,0\}\\
	&=&f_{x}(y) - y.
\end{eqnarray*}
From this it follows that
$$
	\prob\left(\bm{m}_{R}(v) = \bm{m}_{R}(w)\,|\,\mathcal{V}' \right) \leq n^{-\frac{1}{2}f_{x}(y) + O(\log^{-1}n)},
$$
which establishes \eqref{eq:to_show} for this case. The proof of the theorem is finished.
\end{proof}

\end{document}
